\documentclass[11pt,a4paper]{article}
\usepackage[margin=25mm]{geometry}
\usepackage{amsmath,amssymb,amsthm,hyperref}
\setlength{\emergencystretch}{2em}
\newtheorem{lemma}{Lemma}
\newtheorem{theorem}[lemma]{Theorem}
\newcommand{\gen}[1]{\langle#1\rangle}
\title{One neighbour of a binary block swap}
\author{Human Compression project: a component of the L044 recovery}
\date{}
\begin{document}\maketitle
Let the natural $A_{2m}$ have a string C-marking of rank $m$,
where $m\ge8$ is even. Let $H=G_i$ be a transitive maximal
omission preserving two blocks of size $m$, called its rows.
Assume that every distinguished maximal omission $H_j=G_{i,j}$
is intransitive. Choose a distinguished involution $l$ of $H$
interchanging the rows, and suppose precisely one other mark $h$
of $H$ fails to commute with $l$. Write $C$ for the other $m-3$
marks. Throughout, we assume the alternating rank bound
$\lfloor(q-1)/2\rfloor$ in every smaller degree $12\le q<2m$;
the exact acquired finite maxima apply below degree twelve.
We acquire the completed pair-block theorem under this induction,
as well as the complete transitive rank theorem and the extremal
symmetric independent-generator theorem used below.
We also acquire the common contracted-graph reduction, whose
actual-action support statements are recalled below.

\begin{theorem}\label{main}
If $h$ preserves the two rows, this configuration is impossible.
If $h$ interchanges them, choosing $h$ instead of $l$ as the
distinguished block swap gives exactly two noncommuting neighbours
in $H$, with the same omitted label and the same binary partition.
Thus the latter case belongs to the separate two-neighbour branch.
\end{theorem}

The $m$ orbits of $l$ will be called pairs. Contracting them creates
only a graph. The mark $h$ need not act on its vertices as a
permutation. Marks in $C$ do act on the pairs: in coordinates
$(x,\epsilon)$ matching the two rows through $l$, they have the form
\[
                 (x,\epsilon)\longmapsto
                       (x^{\pi},\epsilon+\chi).
                                                        \tag{1}
\]
For the group $\gen C$ the projection to $S_m$ is faithful.
Its kernel lies in $\gen l$, and IP excludes $l$ from $\gen C$.
Consequently $\chi$ is a character of the actual projected group.

The common reduction supplies an eligible selected edge for each
mark other than $l$. Eligibility means that its endpoints lie in
distinct actual $H_j$-orbits. The resulting contracted forest has
two components. The proved local support alternatives say that if
$h$ preserves the rows, it is a double transposition whose contracted
edges are distinct, either disjoint or a path of length two. If it
interchanges the rows, its contracted graph is a triangle on three
pairs and it agrees with $l$ on all other pairs. These assertions
use actual matchings and evenness, not a quotient action of $h$.

\begin{lemma}\label{tree}
Suppose $h$ preserves the rows. Both its contracted edges, together
with one selected edge for every mark in $C$, form a tree $T$ on
all $m$ pairs. Every other contracted edge of a mark in $C$ has
both endpoints among the pairs touched by $h$.
\end{lemma}
\begin{proof}
Let $J$ consist of both $h$-edges. Its graph rank is two. Every
selected $c$-edge is a bridge in the graph consisting of $J$ and
all selected $C$-edges. An alternative path would use no occurrence
of its label; lifting that path by inserting $l$-edges would join
its endpoints inside $H_c$, contradicting eligibility. Adding
$m-3$ bridges gives graph rank $m-1$, so the graph is a connected
tree on $m$ vertices.

Let $e$ be an additional contracted $c$-edge, and consider its unique
path in $T$. That path cannot contain a selected edge of another
$C$-label $d$: the new $c$-edge and the rest of the path would give
a $d$-free bypass, which lifts to $H_d$. Thus its labels are only
$c,h$. If its selected $c$-edge occurs, neither end of the path
can be on that edge, since $c$ induces an involution on the pairs
and its distinct transposition edges are disjoint. The first and
last path edges must therefore both be the two $h$-edges; the
path is $h,c,h$. Otherwise the path uses only $h$. In either case
both endpoints of $e$ belong to the set $Y$ of pairs touched by
$h$. Flips within one pair are omitted from this contracted graph
and remain controlled by (1).
\end{proof}

Every mark of $C$ except a possible neighbour $c$ of $h$ centralizes
both $l,h$. If the two contracted $h$-edges are disjoint, the joint
group has two four-point orbits over $Y$; if they meet, it has one
six-point orbit there. All its other orbits have size two. A
centralizer preserves this union of the larger orbits, so these
marks preserve $Y$. Transitivity and Lemma~\ref{tree} imply that
$c$ exists and its selected edge is the unique edge of $T$ leaving
$Y$. There are no additional edges leaving $Y$ by that lemma.

\begin{lemma}\label{disjoint}
The two contracted $h$-edges cannot be disjoint.
\end{lemma}
\begin{proof}
Here $Y$ consists of four vertices. The unique exit from $Y$
forces the two $h$-components to be joined inside $Y$ in $T$.
There is exactly one joining selected edge, labelled $g\in C$,
and the inside path has labels $h,g,h$. The mark $g$ is different
from $c$, whose selected edge leaves $Y$. Since $g$ centralizes
$l,h$, it interchanges the two four-point joint orbits, and
therefore moves every vertex of $Y$ in projection. In particular
it does not commute with $c$: the latter sends the attachment
vertex outside $Y$, whereas $g$ sends it to another vertex of $Y$.
The two compositions disagree because $c$ has only one exit.

There are $m-4\ge4$ outside vertices. Their tree is connected,
so an edge $d$ is incident with the outside endpoint of $c$.
It also fails to commute with $c$. In projection, $d$ sends that
outside endpoint to a different outside vertex, whereas its action
on the inside endpoint remains within $Y$; applying $c$ there
cannot reach the new outside vertex. Thus $h,g,d$ are three
distinct noncommuting neighbours of $c$, impossible in a string.
\end{proof}

\section*{A path on three pairs}
Name the three pairs $a,u,b$, and choose their row coordinates so
that
\[
                   h=(a_0\ u_0)(u_1\ b_1).            \tag{2}
\]
The graph of $l,h$ on their six points is the path
\[
                 a_1,a_0,u_0,u_1,b_1,b_0.
\]
Every mark $g\in C\setminus\{c\}$ restricts to an automorphism
of this labelled path, hence to its identity or reversal.
In projection on $Y$ these are respectively the identity and
$\tau=(a\ b)$; their row characters are respectively zero and one.
The selected $c$-edge joins a vertex of $Y$ to an outside vertex
$x$. The outside tree is connected.

\begin{lemma}\label{centralpair}
If that vertex is the middle pair $u$, the completed pair-block
branch excludes the configuration.
\end{lemma}
\begin{proof}
Put $z'=(lh)^3$. On the six-point path this is reversal, pairing
$a_0$ with $b_1$, $a_1$ with $b_0$, and $u_0$ with $u_1$.
Outside it, $z'=l$. Thus $z'$ is fixed-point-free, belongs to $H$,
interchanges the binary rows and centralizes $l,h$. All marks other
than $c$ centralize it. In projection $c$ has the edge $(u\ x)$
and possibly $(a\ b)$, its only possible extra edge by matching
and Lemma~\ref{tree}. Either permutation preserves the pairs of
$z'$, with either row character, and commutes with $z'$. Hence
$z'$ is central in $H$.

We justify maximality of this new pair system. In coordinates
matching the rows through $z'$, any permutation preserving their
binary partition and centralizing $z'$ lies in
$\operatorname{diag}(S_m)\times\gen{z'}$. Let $P$ be the
transitive image of $H$ on the new pairs. If all $m-1$ projected
marks are independent, Whiston makes $P=S_m$. Otherwise choose a
redundant projected mark $j$. The actual $H_j$ still projects onto
$P$, but is intransitive on the $2m$ points. It has two orbits,
each meeting every new pair once, so its projection is faithful.
These two orbits are sections of the new pairs, and need not
be the original binary rows.
It supplies a genuine transitive string C-marking of $P$ of rank
$m-2$. The complete acquired L047 theorem makes $P=S_m$ or $A_m$:
for $m\ge10$, $m-2>m/2+1$, and the complete degree-eight proper
transitive boundary has no rank-six case. Thus the new pair system
is maximal. The completed pair-block theorem for the parent
$A_{2m}$ now excludes it, with its stated induction hypotheses.
\end{proof}

We may therefore attach at an outer pair, say $a$, so that the
selected edge is $(a\ x)$. A mark $g\ne c$ with extra edge
$\tau=(a\ b)$ cannot commute with $c$, by its unique exit from
$Y$. Since $c$ already has $h$ as a noncommuting neighbour, at most
one such $g$ exists. The selected edge at the outside endpoint $x$
also fails to commute with $c$, so that possible $g$ must label
that first outside edge. The outside tree has degree at most two:
three incident edges would give three pairwise noncommuting
projected transpositions and a triangle in the string diagram.
Its attachment vertex $x$ has degree one, because two incident
outside edges would give $c$ two noncommuting neighbours in
addition to $h$. Thus this tree is a path starting at $x$.
Write its first edges as $(x\ y),(y\ z),\ldots$. There is a
second edge because the outside set has $m-3\ge5$ vertices.

\begin{lemma}\label{tailtwist}
No $g\ne c$ has the extra edge $(a\ b)$.
\end{lemma}
\begin{proof}
If it did, its projected mark would be $(a\ b)(x\ y)$.
The projection of $c$ could not also have $(b\ u)$: then the
whole projected $C$ graph would be connected, making
$H_h=\gen{l,C}$ transitive, contrary to hypothesis. Thus the
projected first marks are
\[
                  (a\ x),\quad(a\ b)(x\ y),\quad(y\ z),\ldots.
\]
They generate the natural $S_{m-1}$ fixing $u$. The first two
form a transitive dihedral core on four indices; conjugating the
first attached transposition by that core supplies all star
transpositions on five indices. Subsequent path transpositions
adjoin one index at a time. The second tail mark $(y\ z)$ has
row character zero, since its restriction to the six-point path
is the identity. A character of $S_{m-1}$ vanishing on a transposition
is trivial, as transpositions are conjugate and generate the group.
But the character of $g$ is one, because it reverses the six-point
path. This contradiction proves the claim.
\end{proof}

All marks of $C$ except $c$ are consequently ordinary doubled path
transpositions with character zero. The projection of $c$ is
$(a\ x)$, possibly multiplied by $(b\ u)$. If the extra factor
is absent, the projected group is the natural $S_{m-2}$ on $a$
and the outside indices. Its character is trivial, since one of
its transposition generators has character zero. This is the
canonical case treated below. We next exclude the extra factor.

\begin{lemma}\label{localip}
Let $\bar c$ be the doubled permutation $(a\ x)(b\ u)$ and
let $d$ be the doubled transposition $(x\ y)$. Neither
$c=\bar c$ nor $c=\bar c l$ is compatible with the intersection
property of the actual marks $l,h,c,d$.
\end{lemma}
\begin{proof}
Write $c_0,c_1$ for the restrictions of $\bar c$ to the two rows.
First take $c=\bar c$. The element $v=(hc)^2$ is the pure row-zero
double transposition $(a\ u)(b\ x)$. Its conjugate by $lhl$
is the different double transposition $(a\ b)(u\ x)$ on that
row, and their product is $c_0$. Thus $\gen{l,h,c}$ contains
$c_0$ and $c_1=c c_0$.

On the other hand, $e=(cd)^3$ is the doubled $(b\ u)$, so $he$
is a pure row-zero three-cycle on $\{a,b,u\}$. It and its
$c$-conjugate generate $A_4$ on $\{a,b,u,x\}$. Indeed they
are even, their supports make their action transitive on four
points, and their generated order is divisible by four and three;
it must equal twelve. Hence $\gen{h,c,d}$ also contains $c_0$
and $c_1$. But $\gen{h,c}$ is dihedral of order eight, with a
faithful action on row zero: there its two generators have product
of order four. It cannot contain the nonidentity $c_1$ supported
only on row one. This contradicts
\[
            \gen{l,h,c}\cap\gen{h,c,d}=\gen{h,c}.
\]

Now take $c=\bar c l$. This time $e=(cd)^3$ is the doubled
$(b\ u)$ multiplied by $l$. Put $q=(he)^2$. It is the diagonal
copy of the preceding three-cycle on $\{a,b,u\}$.
The elements $q,q^c$ generate diagonal $A_4$ on
$\{a,b,u,x\}$. Together with $d$ they generate diagonal $S_5$
on those four indices and $y$: conjugating $d$ by the transitive
$A_4$ gives all star transpositions to $y$. In particular
$\gen{h,c,d}$ contains $\bar c$, and hence contains
$l=\bar c c$. This contradicts IP between the singleton mark
$l$ and the disjoint marked set $\{h,c,d\}$. Both exclusions
are uniform permutation identities; no finite enumeration is
used as a premise.
\end{proof}

\section*{The canonical case and the missing generator}
Only the following actual marking of $H$ remains: $l$ flips all
pairs, $h$ is (2), and $C$ is the doubled Coxeter path generating
$S_O$, where $O$ consists of $a$ and all the outside indices.
Its noncommutation diagram is the connected path
$l,h,c,d,\ldots$. Therefore the omitted label $i$ is an endpoint
of the full string; removing an interior label would leave two
nonempty commuting components. There are two orientations for
the omitted even involution $\sigma$.

If it is before $l$, it centralizes $J=\gen{h,C}$. The actual
$J$-orbits have sizes $m-1,m-2,2,1$, on
\[
 O_0\cup\{u_0\},\qquad O_1,\qquad
                \{u_1,b_1\},\qquad\{b_0\}.
\]
These sizes are distinct. The induced groups on the first two
orbits are the full natural $S_{m-1},S_{m-2}$, whose centralizers
are trivial. Thus $\sigma$ fixes those orbits pointwise and fixes
the singleton. Its only possible nontrivial action is the
transposition $(u_1\ b_1)$, which is odd. No nonidentity even
$\sigma$ is possible.

If it is after the tail, it centralizes $K$ generated by $l,h$
and all but the last mark of $C$. Let $t$ be the final index.
The actual $K$-orbits are the pair over $t$ and its complement
of size $2(m-1)$. Transitivity on that complement follows from
the connected contracted path and its $l$-edges. On it, the
remaining $C$ marks generate diagonal $S_{O'}$ with
$O'=O\setminus\{t\}$. Conjugating $h$ by this group gives
\[
                    h_x=(x_0\ u_0)(u_1\ b_1)
                         \quad(x\in O').
\]
The products $h_xh_y$ generate the pure alternating group on
$O'_0\cup\{u_0\}$, of size $m-2$. A permutation centralizing
$K$ preserves this support and centralizes its natural alternating
action, whose centralizer is trivial. It fixes a point of the
large $K$-orbit, hence fixes that whole orbit by transitivity.
Thus $\sigma$ can only flip the isolated pair over $t$, again
an odd transposition. This excludes the second orientation and
proves the row-preserving part of Theorem~\ref{main}.

Finally, if $h$ interchanges the binary rows, the acquired support
lemma gives one six-point orbit of $\gen{l,h}$ and all other
orbits of size two. If every mark in $C$ commuted with $h$, it
would centralize this group and preserve that unique six-point
orbit, contradicting transitivity of $H$. Hence another neighbour
$c$ of $h$ belongs to $C$. String relations allow no others.
Choosing $h$ as the distinguished block swap therefore makes its
noncommuting remainder precisely $\{l,c\}$. This is the exact
same-partition reduction to the two-neighbour case asserted by
the theorem.

\paragraph{Attribution and scope.}
The target is the one-neighbour part of source Proposition~3.5,
whose graph strategy and displayed canonical path remain attributed
to Cameron--Fernandes--Leemans--Mixer. The safe contracted support
and bridge arguments were developed by the root investigator;
the present continuation supplies the alternative-pair maximality,
character obstruction, explicit local IP contradictions and actual
endpoint centralizers. The row-swapping case remains a reduction
to the separate two-neighbour case. No whole-source recovery is
claimed here.
\end{document}
